%!TEX = ../../root
\subsection{Problem 3}
\subsubsection{(\rom{1})}

\begin{proof}
    Suppose that P(1) is true and 
    \begin{equation*}
        \paren{\forall n \in \N} \bracks{\paren{P(1) \wedge P(2) \wedge \cdots \wedge P(n)} \implies P(n + 1)}. 
    \end{equation*}

    Let 
    \begin{equation*}
        T(n) := P(1) \wedge P(2) \wedge \cdots \wedge P(n). 
    \end{equation*}

    Since $P(1)$ is true, it implies that $T(1)$ is also true. 
    We also assume $T(n)$ is true. 

    By the assumption of strong induction, we obtain
    \begin{equation*}
       T(n) \implies P(n + 1), 
    \end{equation*} 
    
    and hence, since $T(n)$ is already assumed to be true, 
    \begin{equation*}
        T(n) \wedge P(n + 1). 
    \end{equation*}

    By the definition of $T(n + 1)$, it follows that 
    \begin{equation*}
        T(n) \quad \implies \quad T(n + 1). 
    \end{equation*}

    Following the Principle of Mathematical Induction, this yields
    \begin{equation*}
        \paren{\forall n \in \N} T(n) \quad \implies \quad \paren{\forall n \in \N} P(n). 
    \end{equation*} 

    Thus, Principle of Strong Mathematical Induction is correct. 

\end{proof}

\subsubsection{(\rom{2})}

\begin{proof}
    We define$S$ is a set such that 
    \begin{equation*}
        S := \set{n \in \N \cup \set{0} :  p_n(x) \text{ is a degree $n$ polynomial and the coefficient of } x^n \text{ in } p_n(x) \text{ is }1}
    \end{equation*}

    We know $0 \in S$ is true since $p_0(x) = 1x^0 = 1 \cdot 1 = 1$ is a degree 0 polynominal and the coefficient of $x^0$ is $1$. 
    
    We also know that $1 \in S$ is true since $p_1(x) = 1 \cdot x$ is a degree 1 polynominal and the coefficient of $x$ is 1. 

    It is suffice to show that 
    \begin{equation*}
        (n - 1 \in S) \wedge (n \in S) \implies (n + 1 \in S). 
    \end{equation*}

    Suppose $n - 1 \in S$ and $n \in S$ are true. We define 
    \begin{equation*}
        p_{n - 1}(x) = x^{n - 1} + q(x) \quad \text{and} \quad p_{n} = x^n + r(x), 
    \end{equation*} 

    where 
    \begin{equation*}
        \deg r \leq n - 1 \quad \text{and} \quad \deg q \leq n - 2. 
    \end{equation*}

    Then, by the given statement 
    \begin{align*}
        p_{n + 1}(x) &= x \cdot p_{n}(x) - p_{n - 1}(x) \\
        &= x \cdot \paren{x^{n} + r(x)} - (x^{n - 1} + q(x)) \\
        &= x^{n + 1} + xr(x) - x^{n - 1} - q(x). 
    \end{align*}

    Since the term $xr(x) - x^{n - 1} - q(x)$ has at msot degree $n$, the term $x^{n + 1}$ cannot be canceled. 
    Therefore, $p_{n + 1}(x)$ is a degree $n + 1$ polynominal and the coefficient of $x^{n + 1}$ is $1$ 
    
    Thus, by the Principle of Complete Mathematical Induction, 
    \begin{equation*}
        \forall n \in \N \cup \set{0}, n \in S, 
    \end{equation*}

    which implies that the given statement is true. 

\end{proof}

\subsubsection{\rom{3}}
\begin{proof}
    Let $0 < \abs{t} < 1$ and define statement 
    \begin{equation*}
        T(n): p_n(t + t^{-1}) = \frac{t^{n + 1} - t^{-n-1}}{t - t^{-1}}. 
    \end{equation*}

    We know that $T(0)$ is true since 
    \begin{equation*}
        p_0(t + t^{-1}) = \frac{t^{0 + 1} - t^{0 - 1}}{t - t^{-1}} = \frac{t - t^{-1}}{t - t^{-1}} = 1.  
    \end{equation*}

    We also know that $T(1)$ is true since 
    \begin{equation*}
        p_1(t + t^{-1}) = \frac{t^{1 + 1} - t^{-1 - 1}}{t - t^{-1}} = \frac{t^2 - t^{-2}}{t - t^{-1}} = \frac{\paren{t^1 - t^{-1}} \paren{t^1 + t^{-1}}}{t - t^{-1}} =  t + t^{-1}.  
    \end{equation*}

    Suppose that $T(n - 1)$ and $T(n)$ are true. 
    It suffices to show that 
    \begin{equation*}
        \paren{T(n - 1) \wedge T(n)} \implies T(n + 1). 
    \end{equation*}

    Following from the given statement, 
    \begin{align*}
        p_{n + 1}(t + t^{-1}) &= \paren{t + t^{-1}} p_n(t + t^{-1}) - p_{n - 1}(t + t^{-1}) \\
        &= \paren{t + t^{-1}} \cdot  \frac{t^{n + 1} - t^{-n - 1}}{t - t^{-1}} - \frac{t^{n - 1 + 1} - t^{-(n - 1) - 1}}{t - t^{-1}} \\
        &= \paren{t + t^{-1}} \cdot  \frac{t^{n + 1} - t^{-n - 1}}{t - t^{-1}} - \frac{t^{n} - t^{-n}}{t - t^{-1}} \\
        &= \frac{t^{n + 2} - t^{-n + 2}}{t - t^{-1}} \\
        &= \frac{t^{(n + 1) + 1} - t^{-(n + 1) - 1}}{t - t^{-1}}
    \end{align*}

    This implies $T(n + 1)$ is true. 

    By the Principle of Complete Mathematical Induction, we get 
    \begin{equation*}
        (\forall n \in \N\cup\set{0})T(n). 
    \end{equation*}

    Thus, 
    \begin{equation*}
        (\forall n \in \N\cup\set{0})\bracks{p_n(t + t^{-1}) = \frac{t^{n + 1} - t^{-n-1}}{t - t^{-1}}}.
    \end{equation*}
\end{proof}

\subsubsection{\rom{4}}
\begin{proof}
    Given function 
    \begin{equation*}
        g: D \to \R, g(t) := t + t^{-1}
    \end{equation*}
    where 
    \begin{equation*}
        D := \set{x \in \R : 0 < \abs{x } < 1}. 
    \end{equation*}

    We also define 
    \begin{equation*}
        R := \set{g(t) : t \in D}. 
    \end{equation*}

    In order to show that 
    \begin{equation*}
        R = \paren{-\infty, -2} \cup \paren{2, \infty}, 
    \end{equation*}

    we need to prove both 
    \begin{align}
         &R \subseteq \paren{-\infty, -2} \cup \paren{2, \infty} \label{3.4.1}\\ 
         &R \supseteq \paren{-\infty, -2} \cup \paren{2, \infty} \label{3.4.2}
    \end{align}

    Let us start with the proof of \eqref{3.4.1}

    For all $y \in R$, we have $x \in D$ such that 
    \begin{equation*}
        y = x + \frac{1}{x}. 
    \end{equation*}

    Since $x \neq 0$, we get 
    \begin{equation*}
         y = x + \frac{1}{x} \iff xy = x^2 + 1 \iff 0 = x^2 - yx + 1. 
    \end{equation*}

    By the definition of the set $D$, element x must be in real number set. So the equation
    \begin{equation*}
        0 = x^2 - yx + 1
    \end{equation*}

    must have real roots. Therefore, 
    \begin{equation*}
        (-y)^2 - 4*(1)*(1) \geq 0 \implies y^2 - 4 \geq 0 \implies y^2 \geq 4 \implies \abs{y} \geq 2. 
    \end{equation*}

    However, $\abs{y} \neq 2$ since 
    \begin{equation*}
        x = \frac{-(-y) \pm \sqrt{(-y)^2 - 4 \cdot a \cdot 1}}{2 \cdot 1} = \frac{-(2) \pm \sqrt{(-2)^2 - 4 \cdot a \cdot 1}}{2 \cdot 1} = \frac{-2}{2} = -1
    \end{equation*}

    and 
     \begin{equation*}
        x = \frac{-(-y) \pm \sqrt{(-y)^2 - 4 \cdot a \cdot 1}}{2 \cdot 1} = \frac{-(-2) \pm \sqrt{(2)^2 - 4 \cdot a \cdot 1}}{2 \cdot 1} = \frac{2}{2} = 1
    \end{equation*}

    which makes $x \notin D$. 
    Thus, $\abs{y} > 2$, which also implies $R \subseteq \paren{-\infty, -2} \cup \paren{2, \infty} \label{3.4.1}$.
\end{proof}